% !TeX spellcheck = en-US
% !TeX encoding = utf8
% !TeX program = pdflatex
% !BIB program = biber
% -*- coding:utf-8 mod:LaTeX -*-


% vv  scroll down to line 200 for content  vv


\let\ifdeutsch\iffalse
\let\ifenglisch\iftrue
\input{pre-documentclass}
\documentclass[
  % fontsize=11pt is the standard
  a4paper,  % Standard format - only KOMAScript uses paper=a4 - https://tex.stackexchange.com/a/61044/9075
  twoside,  % we are optimizing for both screen and two-sided printing. So the page numbers will jump, but the content is configured to stay in the middle (by using the geometry package)
  bibliography=totoc,
  %               idxtotoc,   %Index ins Inhaltsverzeichnis
  %               liststotoc, %List of X ins Inhaltsverzeichnis, mit liststotocnumbered werden die Abbildungsverzeichnisse nummeriert
  headsepline,
  cleardoublepage=empty,
  parskip=half,
  %               draft    % um zu sehen, wo noch nachgebessert werden muss - wichtig, da Bindungskorrektur mit drin
  draft=false
]{scrbook}
\input{config}


\usepackage[
  title={Comparative Analysis of zk-SNARK Instantiations for Ensuring Ballot Validity},
  author={Felix R\"ohr.},
  type=bachelor,
  institute=SEC, % or other institute names - or just a plain string using {Demo\\Demo...}
  course={Informatik},
  examiner={Prof.\ Dr.\ Ralf K\"usters},
  supervisor={Nicolas Huber,\ M.Sc.},
  startdate={November 8, 2024},
  enddate={May 8, 2025}
]{scientific-thesis-cover}

% Hier stehen alle Abkürzungen
% TEST
\newacronym{er}{ER}{error rate}
\newacronym{fr}{FR}{Fehlerrate}
\newacronym[plural={RDBMS},shortplural={RDBMS}]{rdbms}{RDBMS}{Relational Database Management System}

% ACTUAL ACRONYMS
\newacronym[plural={Zero Knowledge Proofs}, shortplural={ZKPs}]{zkp}{ZKP}{Zero Knowledge Proof}
\newacronym{eeg}{EEG}{Exponential ElGamal Encryption}
\newacronym{pki}{PKI}{Public Key Infrastructure}

\makeindex

\begin{document}

%tex4ht-Konvertierung verschönern
\iftex4ht
  % tell tex4ht to create pictures also for formulas starting with '$'
  % WARNING: a tex4ht run now takes forever!
  \Configure{$}{\PicMath}{\EndPicMath}{}
  $ % <- syntax highlighting fix for emacs
  \Css{body {text-align:justify;}}

  %conversion of .pdf to .png
  \Configure{graphics*}
  {pdf}
  {\Needs{"convert \csname Gin@base\endcsname.pdf
      \csname Gin@base\endcsname.png"}%
    \Picture[pict]{\csname Gin@base\endcsname.png}%
  }
\fi

%\VerbatimFootnotes %verbatim text in Fußnoten erlauben. Geht normalerweise nicht.

\input{commands}
%\newcommand{\todo}[1]{\textcolor{red}{\textbf{TODO:} #1}}
\pagenumbering{arabic}
\Titelblatt

%Eigener Seitenstil fuer die Kurzfassung und das Inhaltsverzeichnis
\deftriplepagestyle{preamble}{}{}{}{}{}{\pagemark}
%Doku zu deftriplepagestyle: scrguide.pdf
\pagestyle{preamble}
\renewcommand*{\chapterpagestyle}{preamble}



%Kurzfassung / abstract
%auch im Stil vom Inhaltsverzeichnis
\section*{Abstract}

<Short summary of the thesis>

\cleardoublepage

%Solely for German courses of study
\section*{Kurzfassung}

<Kurzfassung der Arbeit>

\cleardoublepage


% BEGIN: Verzeichnisse

\iftex4ht
\else
  \microtypesetup{protrusion=false}
\fi

%%%
% Literaturverzeichnis ins TOC mit aufnehmen, aber nur wenn nichts anderes mehr hilft!
% \addcontentsline{toc}{chapter}{Literaturverzeichnis}
%
% oder zB
%\addcontentsline{toc}{section}{Abkürzungsverzeichnis}
%
%%%

%Produce table of contents
%
%In case you have trouble with headings reaching into the page numbers, enable the following three lines.
%Hint by http://golatex.de/inhaltsverzeichnis-schreibt-ueber-rand-t3106.html
%
%\makeatletter
%\renewcommand{\@pnumwidth}{2em}
%\makeatother
%
\tableofcontents

% Bei einem ungünstigen Seitenumbruch im Inhaltsverzeichnis, kann dieser mit
% \addtocontents{toc}{\protect\newpage}
% an der passenden Stelle im Fließtext erzwungen werden.

\listoffigures
\listoftables

%Wird nur bei Verwendung von der lstlisting-Umgebung mit dem "caption"-Parameter benoetigt
%\lstlistoflistings
%ansonsten:
\ifdeutsch
  \listof{Listing}{Verzeichnis der Listings}
\else
  \listof{Listing}{List of Listings}
\fi

%mittels \newfloat wurde die Algorithmus-Gleitumgebung definiert.
%Mit folgendem Befehl werden alle floats dieses Typs ausgegeben
\ifdeutsch
  \listof{Algorithmus}{Verzeichnis der Algorithmen}
\else
  \listof{Algorithmus}{List of Algorithms}
\fi
%\listofalgorithms %Ist nur für Algorithmen, die mittels \begin{algorithm} umschlossen werden, nötig

% Abkürzungsverzeichnis
\printnoidxglossaries

\iftex4ht
\else
  %Optischen Randausgleich und Grauwertkorrektur wieder aktivieren
  \microtypesetup{protrusion=true}
\fi

% END: Verzeichnisse


% Headline and footline
\renewcommand*{\chapterpagestyle}{scrplain}
\pagestyle{scrheadings}
\pagestyle{scrheadings}
\ihead[]{}
\chead[]{}
\ohead[]{\headmark}
\cfoot[]{}
\ofoot[\usekomafont{pagenumber}\thepage]{\usekomafont{pagenumber}\thepage}
\ifoot[]{}


%% vv  scroll down for content  vv %%































%%%%%%%%%%%%%%%%%%%%%%%%%%%%%%%%%%%%%%%%%%%%%%%%%%%%%%%%%%%%%%%%%%%%%%%%%%%%%%
%
% Main content starts here
%
%%%%%%%%%%%%%%%%%%%%%%%%%%%%%%%%%%%%%%%%%%%%%%%%%%%%%%%%%%%%%%%%%%%%%%%%%%%%%%


\chapter{Introduction}

In today's world, the necessity of secure e-voting schemes is undisputed. 
In many cases, Zero Knowledge Proofs (ZKPs) are used to allow users to prove the validity of their (encrypted) votes without revealing them. 
One class of ZKPs that can be used to prove (essentially) arbitrary valid statements is the class of General Purpose Zero Knowledge proofs (GPZKPs). 
For efficiency reasons, Zero-Knowledge Succinct Non-interactive Arguments of Knowledge (zk-SNARKs) are often used in practice in order to create such (GP)ZKPs.

This thesis starts with \cref{chap:k2}.

We can also typeset \verb|<text>verbatim text</text>|.
Backticks are also rendered correctly: \verb|`words in backticks`|.

\chapter{Voting Systems}
\label{chap:votingsystems}
Voting systems are algorithms that conduct elections. In an election, multiple voters can submit ballots and a voting authority computes the result of the election fromn these ballots. 
To ensure that nobody, including the voting authority, learns anything about the ballots chosen by the voters, the voters only submit the encryption of their ballots.

In a voting system, a voter needs to execute the following steps:
\begin{itemize}
  \item Choose a ballot $b$ from an allowed range of ballots called choice space $C$.
  \item Encrypt the ballot using some (semi-)homomorphic encryption scheme.
  \item Send the encrypted ballot $\tilde{b}$ to the voting authority.
\end{itemize}
The voting authority does the following computations:
\begin{itemize}
  \item Do some computations on the encrypted ballots $\tilde{b}_1,\tilde{b}_2,\dots,\tilde{b}_n$ submitted by the voters using an aggregation function $agg$.
  \item Decrypt the aggregated ballot $\tilde{b}_total$.
  \item Output the final aggregated tally $b_{total}$.
\end{itemize}

Unfortunately, this approach still has some security issues. Specifically, the following two questions are unanswered by the approach:
\begin{enumerate}
  \item How can we convince the voting authority that the submitted encrypted ballots $\tilde{b}_i$ belong to a valid ballots $b_i$ from the choice space $C$ without obtaining any information about the plain ballots $b_i$?
  \item How can we ensure that the voting authority only decrypts the final tally $\tilde{b}_{total}$ and not any of the individual votes $\tilde{b}_i$?
\end{enumerate}

This thesis will focus on the first question. 
We can solve the second question using distributed decryption where multiple trustees choose their own public keys and decryption is only possible by combining all of them. 
Adida et. al. present this solution for Exponential ElGamal encryption %\cite{helios}.

\section{Proving Ballot Validity}
A voter $V_i$ can employ a Zero-Knowledge-Proof (ZKP) to convince the voting authority in a voting system that the encrypted ballot $\tilde{b}_i$ they submitted belongs to a plain ballot $b_i$ from the choice space $C$. 
By using a ZKP, the voter does not disclose any information about the plain ballot $b_i$ to the voting authority.

We can formalize this as follows:
\begin{definition}[Voter]
  Let $C$ be a non-empty choice space, let $\mathcal{S}=(C, \text{Gen}(1^\eta), enc:C\to\tilde{C}, dec:\tilde{C}\to C)$ be an encryption scheme (\todo{Explain/cite}) with $C\subseteq X$.
  Then, a voter $V$ is defined by the following algorithm:
  
\end{definition}

\begin{definition}[Voting system]
  Let $C$ be a non-empty choice space, let $\mathcal{S}=(C, \text{Gen}(1^\eta), enc:C\to\tilde{C}, dec:\tilde{C}\to C)$ be an encryption scheme (\todo{Explain/cite}) with $C\subseteq X$.
  Let $agg:\tilde{C}^*\to\tilde{C}$ be an aggregation function. Then, the voting system $VS=(\mathcal{S}, agg)$
\end{definition}

For now, 

On closer inspection we can split this into two seperate problems:
\begin{enumerate}
  \item How can we convince the voting authority that the voter $i$ knows a plain ballot $b_i$ that belongs to the encrypted one $\tilde{b}_i$?
  \item How can we convince the voting authority that this plain ballot $b_i$ is valid?
\end{enumerate}

To convince the voting authority of the truth of these two statements, we employ Zero-Knowledge-Proofs.

\section{Membership of a Choice Space}

\section{Validity of Encryption}

\chapter{Preliminaries: GPZKPs, SNARKs, Groth16}
\label{chap:k2}

LaTeX hints are provided in \cref{chap:latexhints}.

\blinddocument

\chapter{Formalizing Statements for Groth16}



% \chapter{Proving Ballot Validity using Groth16}

\chapter{Elliptic Curve Cryptography in Circom}
All signals in Circom are represented as elements from a field $\mathbb{F}_p$ with the order of a large, configurable prime $p$. By default:
\begin{align*}
  p = 21888242871839275222246405745257275088548364400416034343698204186575808495617
\end{align*}
This prime field is chosen because it is the scalar field of the BN254 curve. 
The BN254 curve is the basis for the Groth16 proof system.

To define elliptic curve arithmetic in circom, we need to choose a curve which uses the field $\mathbb{F}_p$ as the base field. 
Curves that are constructed in this way are also called \textit{embedded curves}.
One such option is the Curve25519 family. 

\section{Curve25519}
\subsection{The Group Law}
\subsection{Scalar Multiplication}
\subsubsection{Montgomery Ladder}
\subsubsection{y-Recovery}
\section{Exponential ElGamal}

\chapter{Asserting Ballot Validity in Circom}
\section{Single-Vote}
\section{Multi-Vote}
\subsection{Multi-Vote with Rules (MWR)}
\section{Line-Vote}
\section{Borda Election types}
\subsection{Pointlist-Borda}
\subsection{Borda Tournament Style (BTS)}
\section{Condorcet Methods}

\chapter{snarkjs}
\section{Generating the Powers of Tau file}
\section{Proof generation}
\section{Proof verification}

\chapter{Overall benchmarks for Proving Ballot Validity}

\chapter{Related Work}

Describe relevant scientific literature related to your work.

\chapter{Conclusion and Outlook}
\label{chap:zusfas}

\section*{Outlook}

\printbibliography

All links were last followed on March 17, 2018.

\appendix
\input{latexhints-english}

\pagestyle{empty}
\renewcommand*{\chapterpagestyle}{empty}
\Versicherung
\end{document}
